\section{The cap dilation and its linear limit}\label{sec:dilation}

We now let the physical horizons tend to zero. Choose cap shots with
\(T_N\downarrow0\), terminal annular scales satisfying
\(s_N/T_N\to0\), and cutoffs \(b_N\to\infty\). Relabeling a
subsequence in this way does not require that every integer index occur.
Let \(t_{*,N}\in[0,T_N]\) be a time attaining the cap. In the periodic
lift about the origin define
\begin{equation}\label{eq:rescaling}
\begin{gathered}
r_N=t_{*,N}+s_N,\qquad
h_{r_N}(t,x,v)=r_N^{-1}h_N(r_Nt,r_Nx,v),\qquad
C_{r_N}(t,x,v)=r_Nc_N(r_Nt,r_Nx,v),\\
l_N=\frac{s_N}{r_N}\longrightarrow l\in[0,1],\qquad
U_N=\frac{T_N}{r_N}\longrightarrow U\in[1,\infty].
\end{gathered}
\end{equation}
The limits are taken after subselection, using the compactification of
\([0,\infty)\) for the second ratio. We have \(r_N\to0\), and
\(U\geq1\) follows from
\(U_N\geq(1+s_N/T_N)^{-1}\). The prelimit inequality \(U_N\geq1\)
is not needed. The dilated cap event occurs at time \(1-l_N\), and
\begin{equation}\label{dil:cap-bound}
0\leq h_{r_N}(t,x,v)\leq(t+l_N)H(v),
\qquad 0\leq t\leq U_N.
\end{equation}
When convenient we suppress the index on \(r_N\).

\subsection{The limiting mesh and background}

The dilated annular scales are \(\mu_{N,j}=s_j/r_N\), \(j\leq N\).
Their ratio is \(e^\eta\), their smallest member is \(l_N\), and
their largest member tends to infinity. Subselect their logarithmic
phase modulo \(\eta\). The grids then converge on compact
subintervals of \((0,\infty)\) to a geometric mesh \(\mathcal M\)
unbounded above. If \(l>0\), its members are
\(l,e^\eta l,e^{2\eta}l,\ldots\). If \(l=0\), it is a two-sided
geometric mesh with no smallest member. To verify this description,
choose a grid point in \((e^{-\eta},1]\), pass to its limit, and
multiply it by the fixed factors \(e^{j\eta}\). When the smallest
point stays positive its integer distance in the grid from that chosen
point is bounded, so can also be fixed by subselection. Coincident phase
endpoints give the same mesh.

For each \(\mu>0\) and \(\sigma\in\{-1,1\}\), let
\(u_{\mu,\sigma}(t,z)\) be the unique solution of
\[
u=\sigma+\alpha\tanh\frac{\gamma(z-tu)}\mu.
\]
The left side minus the right side has derivative in \(u\) equal to
\(1+t\alpha\gamma S(\gamma(z-tu)/\mu)/\mu>0\); its values have
opposite signs for sufficiently large positive and negative \(u\).
This proves existence and uniqueness. The implicit function theorem
gives smoothness, and
\[
|u_{\mu,\sigma}-\sigma|\leq\alpha,
\qquad
\partial_z u_{\mu,\sigma}
=\frac{\alpha\gamma S(\gamma(z-tu)/\mu)}
       {\mu+t\alpha\gamma S(\gamma(z-tu)/\mu)}\geq0.
\]
On an open mesh annulus \(\mu<|y|<e^\eta\mu\), define the
nonnegative velocity measure
\begin{equation}\label{eq:background}
\begin{split}
B(t,x,\dd p)
&=\sum_{\sigma=\pm1}b_{\mu,\sigma}(t,z)
                         \delta_{u_{\mu,\sigma}(t,z)e_3}(\dd p),\\
b_{\mu,\sigma}(t,z)
&=\frac{\theta K S(\gamma(z-tu_{\mu,\sigma})/\mu)}
 {\mu+t\alpha\gamma S(\gamma(z-tu_{\mu,\sigma})/\mu)}.
\end{split}
\end{equation}
Set \(B=0\) in the hole \(|y|<l\) if \(l>0\), and assign arbitrary
values obeying the same bounds on annular interfaces and the transverse
axis. These exceptional positions have zero Lebesgue measure. In
particular, for every \(k\geq0\),
\begin{equation}\label{dil:background-bound}
\int\langle p\rangle^kB(t,x,\dd p)
\leq\frac{C_k}{t+\mu}\leq\frac{C_k}{t}
\quad(t>0)
\end{equation}
on an active annulus. If \(l>0\), the bound is instead uniformly
\(C_k/(t+l)\) for all \(t\geq0\).

\begin{proposition}\label{prop:blowup-limit}
Every sequence \eqref{eq:rescaling} has a subsequence for which
\(h_{r_N}\rightharpoonup g\) in local spacetime \(L^1\) on
\((0,U)\times\R^3_x\times\R^3_v\), where
\begin{equation}\label{eq:limit-equation}
0\leq g(t,x,v)\leq(t+l)H(v),\qquad
\D g=2P(B,g)-g\nu_B.
\end{equation}
Here \(B\) is \eqref{eq:background}; its gain and loss are defined by
integration against the velocity measure. Both terms on the right are
locally integrable, with a velocity majorant uniform in position on
each compact positive-time interval. Compactly supported velocity
averages of \(h_{r_N}\) converge strongly in local \(L^1_{t,x}\).
The limit has a locally strongly continuous \(L^1\) representative
at positive times and obeys the integrating-factor formula
\eqref{dil:mild} below.

If \(l>0\), pass also to a limit \(\lambda_N\to\lambda_0\) of the
shot parameters. The conclusion and its local continuity extend to
time zero, with
\begin{equation}\label{dil:initial-trace}
g(0,x,v)=\lambda_0 l\,\chi(x/l,v).
\end{equation}
\end{proposition}

\subsection{Cold convergence: survival, Jacobian, and stripe averaging}

We first identify the cold measure independently of the hot limit.
Fix a compact time-space region
\[
a\leq t\leq A<U,\qquad |x|\leq R,
\]
where \(a>0\). Since \(U_N\to U\) and \(A<U\), eventually
\(r_NA<T_N\). The cap therefore holds on the entire initial
physical interval \([0,r_NA]\), as required for the survival
estimate below.
By Proposition~\ref{prop:cold-control}, the dilated
cold velocity moments are bounded there, and indeed uniformly in
position, by
\begin{equation}\label{dil:cold-moment-bound}
\int\langle p\rangle^k C_{r_N}(t,x,p)\dd p
\leq C_k\left(\frac1{t+l_N}+1\right).
\end{equation}
The escaped tails have the still smaller bound
\(C_M r_N^{M+1}t^Me^{-4|p|^2}\) for any fixed \(M>0\).
Uniformly decreasing velocity tails hold under these moment bounds.

First restrict to annuli with \(\mu_{N,j}\) in a fixed compact
interval \([\varepsilon,R_1]\subset(0,\infty)\). There are only
finitely many of them, uniformly in \(N\). Apply
Corollary~\ref{jet:survival} at physical time \(r_Nt\) to a color
of scale \(s=r_N\mu_{N,j}\). Its optical depth is bounded,
uniformly for \(0\le t\le A\), by
\(C_A(r_Nd_s/s+r_N^2+r_Ns_N)+C_{A,M}r_N^{M+1}\). Here
\[
\frac{r_Nd_s}{s}=d_0r_N^{10}\mu_{N,j}^{9},\qquad
r_N^2+r_Ns_N\leq2r_N^2,
\]
so the corollary gives the same-color comparison
\[
e^{-\epsilon_N}c_{i,\mathrm{free}}\leq c_i
\leq c_{i,\mathrm{free}}+C(d_s/s)m_i,
\qquad \epsilon_N\longrightarrow0
\]
uniformly on these scales. The lower bound covers the full free core
because its velocity parameter lies in the compact support of \(J\).
After multiplication by \(r_N\), the remainder has velocity mass
\begin{equation}\label{dil:remainder-power}
r_N O(d_s/s^2)=O(d_0r_N^9\mu_{N,j}^8).
\end{equation}
The core moment bound and this estimate show that the dilated actual
color and its free core differ by
\(o(1)\) in velocity total variation, uniformly on compact regions
of these annuli. The same holds with each fixed polynomial velocity
weight, using compact support of the free-core parameter and the
Gaussian remainder.

We compute the free core exactly. Fix one sign and one such annulus,
write \(\mu_N=s/r_N\), and use the velocity parameters
\[
\bar p=d_s\bar w,\qquad
p_3=\sigma+\alpha\tanh\frac{\gamma(z-tp_3)}{\mu_N}+d_sw_3.
\]
The axial map depends only on \(w_3\). Its derivative is
\[
\frac{\partial p_3}{\partial w_3}
=\frac{d_s}
 {1+t\alpha\gamma S(\gamma(z-tp_3)/\mu_N)/\mu_N}.
\]
Thus Equation~\eqref{eq:core-jacobian}, after dilation, says that the
velocity integral of the free core against a test \(\psi\) is
\begin{equation}\label{dil:free-core-integral}
\begin{split}
\int r_N c_{i,\mathrm{free}}(r_Nt,r_Nx,p)\psi(p)\dd p
={}&\int J(w)
 \frac{K S(\gamma(z-tp_3)/\mu_N)}
      {\mu_N+t\alpha\gamma S(\gamma(z-tp_3)/\mu_N)}
 \psi(d_s\bar w,p_3)\\
&\quad\cdot\ind_E(r_Ny-r_Ntd_s\bar w)
 \ind_{\{|r_N(z-tp_3)|<1/8\}}\dd w.
\end{split}
\end{equation}
On the compact regions under consideration the last axial indicator
equals one for all sufficiently large \(N\). Since the implicit
velocity map has derivative at least one in \(p_3\),
\[
|p_3-u_{\mu_N,\sigma}(t,z)|\leq d_s|w_3|.
\]
After the mesh subselection, \(\mu_N\to\mu>0\), so the velocity
test and the smooth amplitude in \eqref{dil:free-core-integral}
converge uniformly for \(w\in\supp J\) to
\[
\psi(u_{\mu,\sigma}e_3)
\frac{K S(\gamma(z-tu_{\mu,\sigma})/\mu)}
 {\mu+t\alpha\gamma S(\gamma(z-tu_{\mu,\sigma})/\mu)}.
\]

It remains to average the transverse indicator. In dilated coordinates
its period is
\[
\frac{s^2}{r_N}=r_N\mu_N^2\longrightarrow0.
\]
The transverse shift of its argument is \(td_s\bar w\), whose ratio
to that period is at most
\[
\frac{Ad_s}{r_N\mu_N^2}=A d_0 r_N^9\mu_N^8\longrightarrow0.
\]
For the indicator of the fixed periodic stripe pattern, a translation
by this vanishing relative amount changes its integral on a bounded
set by \(o(1)\), uniformly in its phase and in bounded \(t,w\).
One can verify this by summing the lengths of the small intervals
around stripe endpoints. The annular margins have dilated width
\(O(s^2/r_N)=O(r_N\mu_N^2)\), so their area also tends to zero on
the annulus. The outer annular radii converge; their boundaries can
likewise be removed in sets of arbitrarily small area.

Finally the periodic stripe indicators converge weak-* to their mean
\(\theta\). Explicitly, if \(q\) is the one-period indicator and
\(q-\theta\) has bounded periodic primitive \(Q\), then for a smooth
compactly supported one-dimensional test \(f\), period \(p\), and
arbitrary phase \(\zeta\),
\[
\left|\int f(y_1)(q(y_1/p+\zeta)-\theta)\dd y_1\right|
\leq p\|Q\|_\infty\|f'\|_1.
\]
Fubini's theorem gives the same assertion with the other coordinates
and parameters present. Approximation in \(L^1\) extends it to the
annular restrictions. Integrating \(J\), whose mass is one, in
\eqref{dil:free-core-integral} produces exactly the coefficient
\(b_{\mu,\sigma}\) in \eqref{eq:background}.

These arguments identify the cold limit on all compact positive
scale annuli. To remove the lower scale restriction, use
\eqref{dil:cold-moment-bound}: on the slab \(t\geq a\), the integral
over \(|y|<C\varepsilon\), with bounded \(z\), is at most
\(C_{a,A,R,k}\varepsilon^2\). The same estimate holds for the
candidate \(B\). Scales larger than a fixed multiple of \(R\) have
no main color in the chosen spatial compact set, and their escaped
tails vanish. We have proved that for every continuous compactly
supported velocity test \(\psi\),
\begin{equation}\label{dil:cold-weakstar}
\int C_{r_N}(t,x,p)\psi(p)\dd p
\stackrel{*}{\rightharpoonup}
\int\psi(p)B(t,x,\dd p)
\quad\hbox{in local }L^\infty_{t,x}
\end{equation}
on \((0,U)\times\R^3\). Indeed the computation first gives weak
convergence against smooth compact time-space tests, and the common
\(L^\infty\) bound extends it to all \(L^1\) tests. Polynomial
velocity tails can also be removed uniformly. If \(l>0\), choose
\(\varepsilon<l/2\); all the bounds used above hold down to time
zero, so \eqref{dil:cold-weakstar} does too.

\subsection{The hot equation and the size of every error}

Under \eqref{eq:rescaling}, the derivative of \(r_N^{-1}h_N(r_Nt,r_Nx,v)\)
is the old transport derivative of \(h_N\). Therefore the rescaled
equation is exactly
\begin{equation}\label{dil:prelimit-equation}
\D h_{r_N}
=2P^{b_N}(C_{r_N},h_{r_N})
 -h_{r_N}\nu^{b_N}_{C_{r_N}}+E_N,
\end{equation}
where
\[
\begin{split}
E_N={}&r_N^2\bigl[P^{b_N}(h_{r_N},h_{r_N})
                  -h_{r_N}\nu^{b_N}_{h_{r_N}}\bigr]+S_N,\\
S_N(t,x,v)={}&\sum_{i\ne j}
 P^{b_N}(c_{i,N},c_{j,N})(r_Nt,r_Nx,v).
\end{split}
\]
There is no additional factor in \(S_N\). Its smallness comes from
the physical mixed-cold source estimate: for every sufficiently large
fixed \(M\),
\begin{equation}\label{dil:error-bounds}
\begin{split}
0\leq S_N(t,x,v)&\leq C_M r_N^{M-1}t^{M-1}H(v),\\
r_N^2P^{b_N}(h_{r_N},h_{r_N})
&\leq C r_N^2(t+l_N)^2H(v),\\
r_N^2h_{r_N}\nu^{b_N}_{h_{r_N}}
&\leq C r_N^2(t+l_N)^2\langle v\rangle H(v).
\end{split}
\end{equation}
The second line uses Lemma~\ref{lem:hot-bound}; the last line uses
the first velocity moment of \(H\). Thus on every fixed compact
time interval, \(E_N\to0\) in local spacetime \(L^1\), with an
error bounded by a constant times \(r_N^2+r_N^{M-1}\) after
integration over a fixed spatial compact set. Choose \(M>1\).

The main cold sheets and Lemma~\ref{lem:sheet-bound} also give
\[
P^{b_N}(C_{r_N},H)\leq
 C\left(\frac1{t+l_N}+1\right)H,
\qquad
\nu^{b_N}_{C_{r_N}}(v)\leq
 C\left(\frac1{t+l_N}+1\right)\langle v\rangle.
\]
The escaped cold component is bounded by a vanishing multiple of
\(e^{-4|v|^2}\leq CH\), so its gain is covered by
Lemma~\ref{lem:hot-bound}. Equations \eqref{dil:cap-bound} and
\eqref{dil:prelimit-equation} now imply, on \(a\leq t\leq A<U\),
\begin{equation}\label{dil:transport-majorant}
|\D h_{r_N}(t,x,v)|\leq C_{a,A}\langle v\rangle H(v).
\end{equation}
This is a genuine integrable velocity majorant, uniform in position.
If \(l>0\), the same conclusion holds on \(0\leq t\leq A\), since
\(l_N\geq l/2\) eventually.

The cap gives weak compactness of \(h_{r_N}\) in local \(L^1\),
and \eqref{dil:transport-majorant} permits application of
Lemma~\ref{comp:averaging}. Subselect on an exhaustion of
\((0,U)\times\R^3\). The resulting limit \(g\) satisfies the
bound in \eqref{eq:limit-equation}, and each compactly supported
velocity average converges strongly in local \(L^1_{t,x}\).

\subsection{Passing the cold--hot products}

The dilated cold densities concentrate in velocity, so the
two-density compactness argument of Section~\ref{sec:admissibility}
does not apply to \(C_{r_N}\). The needed product passage uses
\eqref{dil:cold-weakstar} and hot averaging instead.

For compact velocity tests \(\psi,\chi\), put
\[
a_N(t,x)=\int\psi(p)C_{r_N}(t,x,p)\dd p,\qquad
\zeta_N(t,x)=\int\chi(v)h_{r_N}(t,x,v)\dd v.
\]
On every compact positive-time region, \(a_N\) is bounded in
\(L^\infty\) and converges weak-* to
\(a=\int\psi\dd B\), while
\(\zeta_N\to\zeta=\int\chi g\dd v\)
strongly in \(L^1\). For a bounded time-space test \(\varphi\),
\[
\int\varphi(a_N\zeta_N-a\zeta)
=\int\varphi a_N(\zeta_N-\zeta)
 +\int\varphi(a_N-a)\zeta\longrightarrow0.
\]
Finite sums of separated velocity tests approximate continuous tests
uniformly on compact velocity products. The total product mass is
uniformly bounded on compact time-space regions by
\eqref{dil:cold-moment-bound} and the cap. We conclude
\begin{equation}\label{dil:mixed-product}
C_{r_N}(t,x,p)h_{r_N}(t,x,v)\dd p\dd v\dd x\dd t
\rightharpoonup B(t,x,\dd p)g(t,x,v)\dd v\dd x\dd t
\end{equation}
against continuous compact velocity tests, with arbitrary smooth
compact time-space dependence. Uniform moment bounds in both velocities
allow removal of velocity cutoffs in tests with any fixed polynomial
growth for which they are used below.

Here is an explicit weak collision test suitable for this convergence.
For a continuous output test \(\phi\), let
\[
\mathcal K_\phi(p,v)
=\frac{|p-v|}{2}\int_{S^2}
 \phi\left(\frac{p+v}{2}+\frac{|p-v|}{2}\omega\right)\dd\omega.
\]
Lemma~\ref{lem:carleman} gives
\(\int\phi P(c,h)=\iint c(p)h(v)\mathcal K_\phi(p,v)\dd p\dd v\).
The test \(\mathcal K_\phi\) is continuous in both incoming
velocities, including \(p=v\), and is bounded in absolute value by
\(2\pi|p-v|\|\phi\|_\infty\). The loss test is
\(2\pi|p-v|\phi(v)\), with the same properties. Inserting smooth
time-space dependence in \(\phi\) does not change these facts.
On bounded velocity products the factor
\(\ind_{\{|p-v|<b_N\}}\) is eventually one; its complement is
uniformly negligible by the weighted product tails. Applying
\eqref{dil:mixed-product} proves the required gain and loss limits.

The limiting gain is an integrable function, rather than an additional
singular output measure. Indeed each summand of \(B\) is an allowed
zero-width axial sheet in Lemma~\ref{lem:sheet-bound}; hence on
\(a\leq t\leq A\),
\[
0\leq P(B,g)\leq C_{a,A}H,
\qquad 0\leq g\nu_B\leq C_{a,A}\langle v\rangle H.
\]
Equivalently the weak gain measure is dominated by this integrable
function, which also verifies absolute continuity at the velocity
axis. Passing \eqref{dil:prelimit-equation}, using
\eqref{dil:error-bounds}, proves the equation in
\eqref{eq:limit-equation}.

\subsection{A positive lower mesh endpoint and the local mild formula}

When \(l>0\), \eqref{dil:transport-majorant} holds uniformly down to
zero. The rescaled initial hot density, on every fixed spatial compact
set for large \(N\), is exactly
\[
h_{r_N}(0,x,v)=\lambda_N l_N\chi(x/l_N,v).
\]
The other periodic copies have moved outside that compact set.
Subselect \(\lambda_N\to\lambda_0\). Smoothness and compact support
of \(\chi\), and \(l_N\to l>0\), give strong local \(L^1\)
convergence to \eqref{dil:initial-trace}. Passing the weak transport
identity with tests touching time zero is justified by the uniform
source majorant and yields that initial trace for \(g\). The cold
product argument also extends to zero in this case, as already noted.
For \(l=0\) the coefficient can be of order \(1/t\), and the
proposition only asserts the equation on the open positive-time
interval.

On any \([a,A]\subset(0,U)\), the limiting transport source is
locally integrable and has the preceding velocity majorant uniformly
in position. Changing to free coordinates as in
\eqref{comp:path-formula}, first on bounded velocities and enlarged
spatial compact sets, gives a locally strongly continuous \(L^1\)
representative. The uniform integrable velocity tails then remove the
velocity truncation. The cap bound passes to every trace of this
representative. Almost every free path obeys the scalar absolutely
continuous equation
\(\D g+\nu_Bg=2P(B,g)\). Its loss coefficient is locally integrable
along the path, since
\(\nu_B(t,x,v)\leq C_a\langle v\rangle\) for \(t\geq a\).
The integrating factor therefore gives, for \(0<a<t<U\),
\begin{equation}\label{dil:mild}
\begin{split}
g(t,x,v)={}&g(a,x-(t-a)v,v)
 \exp\left[-\int_a^t\nu_B(\tau,x-(t-\tau)v,v)\dd\tau\right]\\
&+2\int_a^t P(B,g)(s,x-(t-s)v,v)
 \exp\left[-\int_s^t\nu_B(\tau,x-(t-\tau)v,v)\dd\tau\right]\dd s.
\end{split}
\end{equation}
This is an almost-everywhere identity, or equivalently an identity in
local \(L^1\) at the displayed times. It holds with \(a=0\) when
\(l>0\). Altering the medium on the null annular interfaces does not
change the identity: free transport preserves phase-space measure, so
Fubini's theorem discards the affected set of paths. These statements
complete the proof of Proposition~\ref{prop:blowup-limit}.
